%! Matrix Computations and A = LU — Unit 2, Lesson 2.3 — Group Activity
\documentclass[10pt]{article}
\usepackage{linalg-article}
\usepackage{linalg-boxes}
\newcommand{\TallMath}[1]{$\displaystyle #1\rule[-1.4em]{0pt}{3.2em}$}
\newcommand{\xx}{\mathbf{x}}
\newcommand{\bb}{\mathbf{b}}
\newcommand{\cc}{\mathbf{c}}

\begin{document}

\pageheader{Unit 2, Lesson 2.3}{Group Activity}
\namepartnerperiod

\vspace{0.10in}

\begin{headlinebox}{blush}
\small\textbf{Factor Once, Solve Fast.} A factorization $A=LU$ \emph{saves} elimination: $U$ is the
upper-triangular result, and $L$ is lower-triangular with $1$s on the diagonal and the multipliers
$\ell$ below. To solve $A\xx=\bb$, sweep twice --- forward $L\cc=\bb$, then back $U\xx=\cc$. Work with
your partner and \emph{show every step}.
\end{headlinebox}

\vspace{0.10in}

\begin{tcolorbox}[colback=white, colframe=black!40, title=\textbf{Tier R --- Remediate},
    fonttitle=\bfseries, arc=1.5mm, left=3mm, right=3mm, top=2mm, bottom=2mm, breakable]
Factor $A=\begin{bmatrix} 2 & 6 \\ 6 & 20 \end{bmatrix}$. The pivot is $2$, so the multiplier is
$\ell=6\div2=3$.
\begin{enumerate}[leftmargin=1.7em, itemsep=9pt]
  \item Eliminate (row$_2 - 3\,$row$_1$) to get $U$. \; $U=\begin{bmatrix} 2 & 6 \\ \blank{0.9cm} & \blank{0.9cm} \end{bmatrix}$ \vspace{0.2cm}
  \item Write $L$: put $1$s on the diagonal and the multiplier below. \; $L=\begin{bmatrix} 1 & 0 \\ \blank{0.9cm} & 1 \end{bmatrix}$ \vspace{0.2cm}
  \item \textbf{Check} $LU=A$ by multiplying. \vspace{1.3cm}
  \item In one sentence: what does the $3$ in $L$ record? \writeline
\end{enumerate}
\end{tcolorbox}

\begin{tcolorbox}[colback=white, colframe=black!40, title=\textbf{Tier A --- Approaching Proficiency},
    fonttitle=\bfseries, arc=1.5mm, left=3mm, right=3mm, top=2mm, bottom=2mm, breakable]
\textbf{One factorization, many orders.} A juice bar blends two mixes. Each liter of \textbf{Mix P}
carries $1$ unit sugar and $3$ units vitamin; each liter of \textbf{Mix Q} carries $2$ sugar and $8$
vitamin. With $p$ liters of P and $q$ of Q the totals are $A\begin{bsmallmatrix} p\\q \end{bsmallmatrix}=\bb$,
where $A=\begin{bmatrix} 1 & 2 \\ 3 & 8 \end{bmatrix}$ (sugar row, vitamin row).
\begin{enumerate}[leftmargin=1.7em, itemsep=8pt]
  \item Factor once: find $\ell$, and write $L$ and $U$. \vspace{1.1cm}
  \item \textbf{Order 1:} sugar $4$, vitamin $14$. Two sweeps ($L\cc=\bb$, then $U\xx=\cc$) --- find $p,q$. \vspace{1.3cm}
  \item \textbf{Order 2:} sugar $7$, vitamin $27$. Same $L$ and $U$ --- find $p,q$. \vspace{1.3cm}
  \item \emph{Interpret:} give each order's liters, and say why keeping $L$ and $U$ beats re-eliminating
        every order. \writeline
\end{enumerate}
\end{tcolorbox}

\begin{tcolorbox}[colback=white, colframe=black!40, title=\textbf{Tier E --- Extension},
    fonttitle=\bfseries, arc=1.5mm, left=3mm, right=3mm, top=2mm, bottom=2mm, breakable]
\textbf{A $3\times3$ factorization.} Factor $A=\begin{bmatrix} 2 & 4 & 2 \\ 4 & 9 & 5 \\ 2 & 7 & 8 \end{bmatrix}$
by elimination, recording all three multipliers $\ell_{21},\ell_{31},\ell_{32}$ into $L$.
\vspace{2.0cm}

\textbf{Then solve} $A\xx=\bb$ for $\bb=\begin{bmatrix} 10 \\ 22 \\ 19 \end{bmatrix}$ using two sweeps
($L\cc=\bb$ forward, then $U\xx=\cc$ back). \emph{Justify:} explain why writing down $L$ cost you no
extra arithmetic beyond the elimination you already did. \writelines{3}
\end{tcolorbox}

\end{document}
